% Auto-generated by generate_narrative_content.py
% Ridge classifier description — addresses reviewer request for classifier details
% Insert in Section III (Architecture) after feature extraction description

\subsection{Classification Stage}

After feature extraction produces an $F$-dimensional feature vector
(where $F = 84 \times D \times P$ for $D$ dilations and $P$ pooling operators),
two additional pipeline stages complete the inference:

\textbf{Scaler Kernel (K2).}
A StandardScaler normalizes each feature $f_i$ to zero mean and unit variance:
$\hat{f}_i = (f_i - \mu_i) / \sigma_i$, where $\mu_i$ and $\sigma_i$ are
precomputed from training data and stored in on-chip BRAM.

\textbf{Classifier Kernel (K3).}
A Ridge regression classifier computes decision scores
$s_c = \mathbf{w}_c^T \hat{\mathbf{f}} + b_c$ for each class $c$,
selecting $\arg\max_c\, s_c$.
Ridge regression is a linear classifier with $L_2$ regularization,
chosen because it matches the feature extraction's
``transform once, classify linearly'' paradigm without iterative optimization.
The coefficient matrix $\mathbf{W} \in \mathbb{R}^{C \times F}$ and
intercept vector $\mathbf{b} \in \mathbb{R}^C$ are precomputed on the host
and transferred via HBM.

K2 and K3 are connected to K1 (feature extraction) via AXI-Stream interfaces,
forming a three-stage pipeline.
For MiniRocket ($F = 840$), K2 and K3 together consume $<$0.1\,ms per sample,
making K1 the dominant bottleneck (85--99\% of total inference time).
